\documentclass{ctexbook}
\usepackage{EC}
\usepackage[
    backend=biber,
    style=chem-acs,
]{biblatex}
\addbibresource{ref.bib}
\begin{document}
\section{氮及其化合物}
\subsection{单质氮}
\subsubsection{氮气~\ce{N2}}
氮气~\ce{N2}是通常情况下氮元素唯一稳定的单质.\paragraph{\ce{N2}的物理性质与结构}
\ce{N2}是一种无色无味无臭的反磁性气体,熔点为$\SI{-210}{\celsius}$,沸点为$\SI{-195.8}{\celsius}$,难溶于水.~\ce{N2}在固相中的结构与~\ce{O2}和~\ce{F2}相似,可以参看对应章节的介绍.
\paragraph{\ce{N2}的化学性质}
\ce{N2}在常温下是相当不活泼的,这主要是因为~\ce{N2}的键解离能为$\SI{945.41}{\kilo\joule\per\mole}$,是非常难以断裂的共价键;同时~\ce{N2}的HOMO-LUMO能级间距大,并且~\ce{N#N}键的电子分布非常均匀,没有极性.~\ce{N2}的稳定性也是目前地球生物存在的必要条件,否则大家可能需要生活在一个充满~\ce{NO}和~\ce{NO2}的大气中.\\
\indent \ce{N2}的反应性随温度升高而增加.在加热条件下~\ce{N2}与~\ce{H2}反应生成~\ce{NH3},与焦炭反应生成\ce{(CN)2}:
\begin{tightcenter}
    \ce{N2 + 3H2 ->[$\Delta$] 2NH3}\\
    \ce{2C + N2 ->[$\Delta$] (CN)2}
\end{tightcenter}
\ce{N2}也能与各种活泼金属(主要是碱土金属,许多过渡金属和~\ce{Al},~\ce{Ge}等主族金属)反应生成氮化物.然而,在碱金属氮化物中,只有~\ce{Li3N}能由~\ce{Li}与~\ce{N2}加热反应得到;~\ce{Na3N}和~\ce{K3N}都是不稳定的化合物,容易分解为单质.这主要是因为~\ce{Li+}半径较小,与高电荷的~\ce{N^3-}形成晶体时有很大的晶格能,而其余碱金属阳离子已经不能很好地稳定~\ce{N^{3-}}.\\
\indent 虽然~\ce{N2}相比其等电子体~\ce{CO}是一个较差的配体(具有更低的HOMO轨道和更高的LUMO轨道),但仍然能形成各种配合物.
\paragraph{\ce{N2}的生产与用途}
大规模制取~\ce{N2}的唯一途径是分馏空气(这和制取~\ce{O2}是一致的),实验室也多采取氮气瓶供给的方式实现.其它的可行的办法主要包括叠氮化物或铵盐的分解:
\begin{tightcenter}
    \ce{2NaN3 ->[$\SI{300}{\celsius}$] 2Na + 3N2}\\
    \ce{NH4NO2(aq) ->[$\Delta$] N2(g) + 2H2O(l)}\\
    \ce{(NH4)2Cr2O7 -> N2 + Cr2O3 + 4H2O}\\
\end{tightcenter}
或氧化剂对~\ce{NH3}的氧化:
\begin{tightcenter}
    \ce{8NH3 + 3Br2 ->T[aq] N2 + 6NH4Br}\\
    \ce{2NH3 + 3CuO -> N2 + 3Cu + 3H2O}
\end{tightcenter}
\subsubsection{新型固相氮单质}
尽管~\ce{N2}极高的稳定性使得其它可能的氮单质都相对不稳定而难以得到,但随着高压化学的发展,人们业已制备或表征了一系列新型固相~\ce{N2}单质,它们都具有复杂层状或网状结构.
\paragraph{cg-N}
2004年, Eremets等人在金刚石对顶砧中将~\ce{N2}压缩至$\SI{115}{\giga\pascal}$并用激光加热到约$\SI{2000}{\kelvin}$以上,首次制得立方晶系的cg-N,这是新型氮单质研究的里程碑\cite{Eremets2004}.\\
\indent cg-N具有古怪的三维网状结构,其中所有~\ce{N}原子的化学环境相同,与周围的三个~\ce{N}原子以~\ce{N-N}单键键合.
\begin{figure}[H]\centering
    \begin{subfigure}[b]{.45\linewidth}
        \centering\includegraphics[scale=0.125]{figure/cg-N-1.png}
        \caption{cg-N的立方晶胞示意图}
    \end{subfigure}
    \begin{subfigure}[b]{.45\linewidth}
        \centering\includegraphics[scale=0.125]{figure/cg-N-2.png}
        \caption{cg-N的晶胞沿$a$轴的投影图}
    \end{subfigure}
    \caption{cg-N的晶体结构\quad$\text{I}\,2_1\,3$\quad$a=b=c=\SI{345.42}{\pico\meter}$}
\end{figure}
\noindent 从任意一个对角线方向看, cg-N都可以视作~\ce{\{N_10\}}环连接成的二维层堆叠成的三维结构(你可以从$a$轴的投影图中找到三组取向的~\ce{\{N_10\}}环,因而判断其属于立方晶系).另一方面,也可以将其视作类似灰硒中的螺旋链被~\ce{N}原子连接的结构.
\begin{figure}[H]\centering
    \begin{subfigure}[b]{.45\linewidth}
        \centering\includegraphics[scale=0.125]{figure/cg-N-3.png}
        \caption{cg-N沿晶胞对角线的投影图}
    \end{subfigure}
    \begin{subfigure}[b]{.45\linewidth}
        \centering\includegraphics[scale=0.125]{figure/cg-N-4.png}
        \caption{cg-N的层状结构示意图}
    \end{subfigure}
    \caption{cg-N的结构解析}
\end{figure}
\noindent 以如此复杂的环为结构基础并且仍然保持立方晶系高对称性的晶体是很少见的.同时, cg-N也具有手性.\\
\indent 2024年, Xu等人优化了cg-N的合成条件,以~\ce{KN3}为原料用等离子体气相沉积法在制备了独立的cg-N, 无需高压或限定在纳米尺度\cite{Xu2024}.他们同时报道常压下cg-N能在约$\SI{760}{\kelvin}$以下稳定存在,高于此温度则发生爆炸性的热分解.
\paragraph{bp-N}
2020年, Laniel等人在约$\SI{140}{\giga\pascal}$下激光加热~\ce{N2}至$\SI{4000}{\kelvin}$得到了具有类似黑磷结构的bp-N\cite{Laniel2020}.
\begin{figure}[H]\centering
    \begin{subfigure}[b]{.45\linewidth}
        \centering\includegraphics[scale=0.125]{figure/bp-N-1.png}
        \caption{bp-N沿晶胞$a$轴投影图}
    \end{subfigure}
    \begin{subfigure}[b]{.45\linewidth}
        \centering\includegraphics[scale=0.125]{figure/bp-N-2.png}
        \caption{bp-N的层状结构示意图}
    \end{subfigure}
    \caption{bp-N的晶体结构\quad$\text{C}\,\text{m}\,\text{c}\,\text{e}$\quad$a=\SI{213.43}{\pico\meter},\,b=\SI{653.42}{\pico\meter},\,c=\SI{285.99}{\pico\meter}$}
\end{figure}
\subsubsection{\ce{N_x}分子}
除去高温高压下合成的固相氮单质外,人们已经在气相中表征或在低温基质中捕获一些新的~\ce{N_x}分子.
\paragraph{\ce{N6}}
2025年, Qian等人报道了~\ce{N6}分子的制备与捕获\cite{Qian2025}:将~\ce{Cl2}或~\ce{Br2}与~\ce{AgN3}在室温下减压反应,将产物收集到$\SI{10}{\kelvin}$的~\ce{Ar}基质上,并进一步在$\SI{77}{\kelvin}$下得到了不依赖基质的纯~\ce{N6}薄膜.这是历史上第一次制备并表征的纯中性多氮分子.\\
\indent \ce{N6}分子是具有$C_{2\text{h}}$构型的链状分子,如下所示:
\begin{figure}[H]\centering
    \includegraphics{figure/N6.eps}
    \caption{\ce{N6}的结构示意图}
\end{figure}
\paragraph{\ce{N4}与~\ce{N3}}
目前只在气相中检测到~\ce{N3}与~\ce{N4}的信号\cite{Cacace2002},没有分离得到相应物质.
\subsection{氮化物}
\subsubsection{简单氮化物}
氮几乎和周期表中的所有元素形成二元化合物(略少于~\ce{O}和~\ce{F}元素).我们可以按照性质把氮化物分成三类:离子型,共价型和金属型.共价型和金属型氮化物的结构通常比较复杂.\\
\indent 金属氮化物的制备方式主要有:
\begin{enumerate}[label=\tbf{\arabic*.}]
    \item 金属与~\ce{N2}或~\ce{NH3}(通常在高温下)反应,例如:
        \begin{tightcenter}
            \ce{3Ca + N2 -> Ca3N2}\qquad\ce{3Mg + 2NH3 -> Mg3N2 + 3H2}
        \end{tightcenter}
    \item 金属氨基化合物的热分解,例如:
        \begin{tightcenter}
            \ce{3Zn(NH2)2 -> Zn3N2 + 4NH3}
        \end{tightcenter}
    \item 在有还原剂存在时还原金属氧化物或卤化物,例如
        \begin{tightcenter}
            \ce{Al2O3 + 3C + N2 -> 2AlN + 3CO}\qquad\ce{2ZrCl4 + N2 + 4H2 -> 2ZrN + 8HCl}
        \end{tightcenter}
    \item 在液氨中进行的反应.
\end{enumerate}
\subsubsection{全氮离子}
自1890年Curtius制备了~\ce{HN3}以来,人们一直致力于对全氮离子的研究,也许主要是它们在高能材料方面的显著优势.除去叠氮阴离子~\ce{[N3]-}以外,近年来已经制备~\ce{[N5]+},~\ce{[N5]-}等一系列结构更复杂的全氮离子.本节将介绍这些离子.
\paragraph{\ce{[N5]+}}
\ce{[N5]+}的发现可以追溯到美国空军从1986年开始的高能量密度物质的研究项目.~1999年,美国空军研究实验室的高级研究员Christe设想了由~\ce{[N2F]+}与~\ce{[N3]-}发生取代反应得到~\ce{[N5]+}的过程,并最终按下面的过程制备得到~\ce{[N5][SbF6]}\cite{Christe1999}:
\begin{tightcenter}
    \ce{cis-N2F2 + SbF5 ->[$\SI{-78}{\celsius}$] [N2F]+[SbF6]-}\\
    \ce{ [N2F]+[SbF6]- + HN3 ->[$\SI{-78}{\celsius}$][\ce{HF}] [N5]+[SbF6]- + HF}
\end{tightcenter}
\ce{N5+}具有折线形结构,所有~\ce{N}原子处在同一平面上:
\begin{figure}[H]\centering
    \includegraphics{figure/N5+.eps}
    \caption{\ce{[N5]+}的折线形结构}
\end{figure}
形成~\ce{N=N=N=N=N}的直线型结构很不利,因为中间的三个~\ce{N}都带形式正电荷.采取折线形结构后,两边的~\ce{-N#N}为了与中心~\ce{N}原子在平面上的$\text{sp}^2$电子共轭而向上弯曲\footnote{仅作一种可能的解释.}.\\
\indent 值得一提的是,~\ce{N5+}可以形成诸如~\ce{[N5]+[P(N3)6]-}和~\ce{[N5]+[B(N3)4]-}等盐.它们也许是含~\ce{N}量较高的化合物之一.
\paragraph{\ce{[N5]-}}
早在2002年,已经有人在气相中检测到环状的~\ce{[N5]-}离子. 2017年, Zhang等人首次合成出室温下稳定的~\ce{[N5]-}盐\cite{Zhang2017}.他们通过芳基五唑的裂解得到了~\ce{[N5]6[H3O]3[NH4]4Cl}:
\begin{figure}[H]\centering
    \includegraphics{figure/N5-.eps}
    \caption{\ce{[N5]-}的合成路线及其结构}
\end{figure}
\ce{[N5]-}的稳定性部分源于其芳香性,它是咪唑阴离子的等电子体.将前述制备的盐~\ce{[N5]6[H3O]3[NH4]4Cl}与~\ce{Co(NO3)2.6H2O}反应可以得到在空气中稳定的橙色晶体~\ce{Co(N5)2(H2O)4.4H2O}\cite{Zhang2017a},其中含有八面体配合物~\ce{Co(N5)2(H2O)4}:
\begin{figure}[H]
    \centering\includegraphics{figure/Co(N5)2(H2O)4.eps}
    \caption{\ce{Co(N5)2(H2O)4}的结构}
\end{figure}
\paragraph{\ce{[N4]^{4-}}}
2019年, Laniel等人在研究~\ce{Mg-N}相的高压化学时,在$\SI{52.2}{\giga\pascal}$和$\SI{1850}{\celsius}$后得到了~\ce{MgN4}和~\ce{Mg2N4}\cite{Laniel2019a}.前者中有~\ce{N}构成的长链阴离子,而后者中含有独立的,具有链状结构的~\ce{[N4]^4-}.键长数据表明~\ce{[N4]^{4-}}中的~\ce{N-N}键有一定多重键性质.
\begin{figure}[H]\centering
    \begin{subfigure}[b]{.55\linewidth}
        \centering\includegraphics[scale=0.125]{figure/Mg2N4.png}
        \caption{\ce{Mg2N4}的晶体结构示意图}
    \end{subfigure}
    \begin{subfigure}[b]{.35\linewidth}
        \centering\includegraphics{figure/[N4]4-.eps}
        \caption{\ce{[N4]^4-}的结构示意图}
    \end{subfigure}
    \caption{\ce{[N4]^4-}的结构}
\end{figure}
\paragraph{\ce{[N6]^{4-}}}
2023年, Laniel等人将~\ce{KN3-N2}混合体系在~$\SI{61}{\giga\pascal}$下加热至$\SI{2000}{\kelvin}$,首次合成了含有~\ce{[N6]^4-}的盐~\ce{K9N56}\cite{Laniel2023},其实际组成为~\ce{[K]18[N6][N5]14[N2]18}.晶体中的~\ce{[N6]^{4-}}呈平面六元环结构,是继~\ce{[N5]-}后又一具有芳香性的全氮阴离子. 2026年, Zhang等人改进了合成方法,得到了仅以~\ce{[N6]^4-}为阴离子的盐~\ce{K4N6}\cite{Zhang2026}.
\subsection{氮的氢化物:氨~\ce{NH3}和铵根离子~\ce{NH4+}}
在所有氮的氢化物中,氨~\ce{NH3}是最重要的一个.~\ce{NH3}是无色的具有特殊刺激性气味的气体;较高浓度的~\ce{NH3}是有毒的.~\ce{NH3}的熔点为$\SI{-77.7}{\celsius}$,沸点为$\SI{-33.4}{\celsius}$.液态~\ce{NH3}具有较低的密度和粘度,但具有很高的介电常数,因此和~\ce{H2O}一样适合作为溶剂.\\
\indent 尽管~\ce{NH3}在凝聚相和~\ce{H2O}一样形成氢键,但不同的是固态~\ce{NH3}中的每个~\ce{N}原子均形成三组氢键,因此这一氢键事实上与冰中具有方向性和饱和性的氢键并不相同,更像是静电相互作用.由于氢键的方向性弱,因此固态氨并不像冰那样出现反常的膨胀,因而融化时实际上沉在液相的底部\footnote{$\SI{198}{\kelvin}$时,液氨的密度为$\SI{0.731}{\gram\per\centi\meter^3}$, $\SI{193}{\kelvin}$时固态氨的密度为$\SI{0.820}{\gram\per\centi\meter^3}$}.
\subsubsection{液氨溶剂}
液氨是人们最熟悉的,研究得最充分的非水离子化体系.在这其中可以发生许多典型的反应.
\paragraph{液氨的活泼金属溶液}
碱金属最引人注意的性质之一是它们易溶于液氮并形成亮蓝色的,具有异乎寻常性质的亚稳态溶液.重碱土金属,~\ce{Eu}和~\ce{Yb}也有类似的性质.\\
\indent 这些溶液具有相似的性质:在浓度较低时,它们都呈现蓝色,而在浓度较高时则显现金属般的赤褐色;随着浓度增大,溶液的摩尔电导率先减小后增大,磁性也由低浓度的顺磁性变为抗磁性,最后又显现微弱的顺磁性.\\
\indent 现在,我们普遍认为其中含有~\ce{[M(NH3)_n]^x+}和分布在空穴中的自由电子~\ce{e-}.这些空穴是~\ce{e-}排开~\ce{NH3}形成的,因此溶液密度明显减小.溶液显现蓝色正是由于自由电子的跃迁吸收.近年来的研究认为可以按照下面的模型描述活泼金属-液氨溶液在不同浓度的性质\cite{Zurek2009}.以~\ce{Li-NH3}溶液为例:\\
\tbf{1.极低浓度}(大约$<10^{-3}\;\si{\mole}\;\%$)\\
\indent \ce{Li}溶于液氨后,其价电子进入溶剂体系.此时主要存在溶剂化锂离子~\ce{[Li(NH3)4]+},溶剂化电子~\ce{e^-(NH3)_m},二者形成的离子对以及中性单体~\ce{Li(NH3)4}:
\begin{tightcenter}
    \ce{[Li(NH3)4]+ + e^-(NH3)_m <=> [Li(NH3)4+.e^-(NH3)_m] <=> Li(NH3)4 + $m$NH3}
\end{tightcenter}
由于此时~\ce{e^-(NH3)_m}等单电子物种浓度很低,相互接近并发生自旋配对的概率很小,因此体系以含一个未配对电子的 $S=1/2$物种为主,溶液呈顺磁性和典型的蓝色.\\
\tbf{2. 中等浓度}(大约$<10^{-3}\sim1\;\si{\mole}\;\%$)\\
\indent 随着浓度升高,不同的单电子物种彼此接近并发生耦合,形成$S=0$的自旋配对的抗磁性缔合物,例如:
\begin{tightcenter}
    \ce{2[e^-(NH3)_m] -> [(e^-)2(NH3)_p]}\qquad\ce{2Li(NH3)4 -> [(Li2)(NH3)_q]}
\end{tightcenter}
随自旋配对程度增加,溶液的顺磁性迅速减弱,磁化率下降.到大约$5\si{\mole}\%$时,实验表明超过$90\%$的电子已经自旋配对.\\
\tbf{3. 高浓度}(大约$1\sim8\;\si{\mole}\;\%$)\\
\indent 在这一浓度范围内,溶液发生向金属态的转变(Transition to the Metallic State, TMS).随着电子耦合物种和~\ce{Li(NH3)4}单元进一步聚集,溶剂化电子逐渐离域,形成类似金属的能带.此时,溶液的摩尔电导率急剧增加,同时颜色由蓝色逐渐转变为青铜色或金色.
\indent 碱金属的液氨溶液不稳定.在有杂质作为催化剂或放置的较久的情况下,这些溶液将分解为对应的氨基盐:
\begin{tightcenter}
    \ce{2Na + 2NH3 -> 2NaNH2 v + H2}
\end{tightcenter}
需要注意的是,~\ce{NaNH2}微溶于液氨,因此放置的太久的~\ce{Na-NH3}溶液发生的现象为:溶液的蓝色逐渐褪去,产生白色沉淀.\\
\indent 由碱金属-胺溶液(与碱金属-液氨溶液性质相似)出发也可以获得以电子或碱金属阴离子的盐. 1974年, Dye等人合成了金色的含~\ce{Na-}的顺磁性物质~\ce{[Na(2,2,2-crypt)]+[Na]-}\cite{Dye1974}:
\begin{tightcenter}
    \ce{2Na + 2,2,2-crypt ->T[\ce{EtNH2(l)}][vap] [Na(2,2,2-crypt)]+[Na]-}
\end{tightcenter}
后来,在相似的体系中又合成了蓝黑色的~\ce{[K(2,2,2-crypt)]+[e^-]}\cite{DeBacker1971}.\\
\indent 碱金属的液氨溶液具有很强的还原性.并且由于还原性是自由电子~\ce{e-}体现的,因此它常常能做到对物质的单电子还原.例如:
\begin{tightcenter}
    \ce{Mn2(CO)10 + 2K ->T[\ce{NH3(l)}] 2K[Mn(CO)5]}\\
    \ce{Fe(CO)5 + 2Na ->T[\ce{NH3(l)}] Na2[Fe(CO)4] + CO}
\end{tightcenter}
在有机化学中,碱金属的液氨溶液也被广泛地运用于还原各种物质,例如将炔烃还原为反式烯烃, Birch还原等等.
\paragraph{液氨中的复分解反应}
大部分时候,我们可以把液氨中的反应与水溶液中的反应进行类比.然而,由于结构与性质上的差异,许多物质在水中的溶解度和在液氨中有明显的不同.典型的微溶物有~\ce{NaF},~\ce{NaCl}等.大多数铵盐都易溶,而~\ce{AgBr}等物质由于可以形成~\ce{Ag(NH3)2+}配离子,也是易溶的.这就造成了某些水溶液中不能进行的反应反而能在液氨中进行:
\begin{tightcenter}
    \ce{Ba(NO3)2 + 2AgBr ->T[\ce{NH3(l)}] BaBr2 v + 2AgNO3}
\end{tightcenter}
以及与难溶氢氧化物和氧化物的形成类似地有
\begin{tightcenter}
    \ce{AgNO3 + KNH2 ->T[\ce{NH3(l)}] AgNH2 v + KNO3}\\
    \ce{3HgI2 + 6KNH2 ->T[\ce{NH3(l)}] Hg3N2 + 6KI + 4NH3}
\end{tightcenter}
\paragraph{液氨中的酸碱反应}
和水一样,液氨也可以发生自耦电离:
\begin{tightcenter}
    \ce{2NH3 <=> NH4+ + NH2-}\quad$K=1.9\times10^{-33}$
\end{tightcenter}
因此体系中可以发生与水溶液中类似的酸碱反应.此外,还有一些物质具有两性:
\begin{tightcenter}
    \ce{K2[Zn(NH2)4] + 2NH4NO3 ->T[\ce{NH3(l)}] Zn(NH2)2 + 2KNO3 + 4NH3}\\
    \ce{Zn(NH2)2 + 2NH4NO3 ->T[\ce{NH3(l)}] [Zn(NH3)4](NO3)2}
\end{tightcenter}
\paragraph{液氨中的氧化还原反应}
不考虑动力学因素时,液氨中不会有比~\ce{N2}更强的氧化剂,也不会有比~\ce{H2}更强的还原剂.这仅仅对应于$\SI{0.04}{\volt}$的电位区间,但由于超电势的存在,实际可存在的氧化剂/还原剂的范围要大得多.\\
\indent 超出上述电位范围的物质仍然会发生缓慢的反应.例如,~\ce{NO3-}在碱性溶液中放出~\ce{N2}:
\begin{tightcenter}
    \ce{3K+ + 3NH2- + 3NO3- ->T[\ce{NH3(l)}] 3KOH v + N2 + 3NO2- + NH3}
\end{tightcenter}
\subsubsection{氨的化学性质}
\ce{NH3}溶于水即生成一水合氨.在水中,~\ce{NH3}是弱碱:
\begin{tightcenter}
    \ce{NH3(aq) + H2O(l) <=> NH4+(aq) + OH-(aq)}\quad$K=1.81\times10^{-5}$
\end{tightcenter}
\indent 氨在空气中很难燃烧,但在催化剂的存在下可以被选择性地氧化为~\ce{NO}.~\ce{NO}接着可以被氧化为~\ce{NO2},这也是工业制~\ce{HNO3}的方法:
\begin{tightcenter}
    \ce{4NH3 + 3O2 ->T[点燃] 2N2 + 6H2O}\\
    \ce{4NH3 + 5O2 ->T[\ce{Pt}][$800\tc$] 4NO + 6H2O}
\end{tightcenter}
此外,利用Andrussow法工业生产~\ce{HCN}也存在~\ce{NH3}的氧化:
\begin{tightcenter}
    \ce{2CH4 + 2NH3 + 3O2 ->[Pt] 2HCN + 6H2O}
\end{tightcenter}
\indent \ce{NH3}在~\ce{F2}中燃烧,形成黄绿色火焰,生成~\ce{NF3}.~\ce{NH3}在~\ce{Cl2}中的燃烧则复杂得多.不过,在正常条件下可以发生反应
\begin{tightcenter}
    \ce{8NH3 + 3Cl2 -> 6NH4Cl + N2}
\end{tightcenter}
这一反应可以用于检测~\ce{Cl2}管道的泄漏:由于~\ce{NH4Cl}为白色固体,因此有~\ce{Cl2}存在的地方会产生明显的白烟.\\
\indent 需要特别注意的是,由于~\ce{Cu}离子与~\ce{NH3}有很好的络合作用,因此有~\ce{O2}存在时~\ce{NH3}将迅速地侵蚀\~ce{Cu}单质.因此,不能用含~\ce{Cu}的管道材料运输~\ce{NH3}.
\subsection{氮的其它氢化物}
除~\ce{NH3}以外,氮还形成联氨~\ce{N2H4},叠氮酸~\ce{HN3}等常见氢化物.
\subsubsection{联氨~\ce{N2H4}}
无水联氨为无色发烟液体,熔点为$\SI{2.0}{\celsius}$,沸点为$\SI{113.5}{\celsius}$.联氨和水的物理性质在很多方面上都比较相似.
\paragraph{联氨的结构}
气相中的联氨呈现歪扭式的结构\cite{Kohata1982}:
\begin{figure}[H]\centering
    \begin{subfigure}{.45\linewidth}\centering
        \includegraphics{figure/N2H4-1.eps}\caption{~\ce{N2H4}的结构}
    \end{subfigure}
    \begin{subfigure}{.45\linewidth}\centering
        \includegraphics{figure/N2H4-2.eps}\caption{~\ce{N2H4}的Newman投影式}
    \end{subfigure}
\end{figure}
不难看出这一结构是平衡了~\ce{N}上孤对电子排斥和~\ce{N}的孤对电子对$\sigma^\ast(\ce{N-H})$轨道的超共轭效应的结果.
\paragraph{联氨的碱性}
由于~\ce{N}的吸电子效应,~\ce{N2H4}在水溶液中的碱性弱于~\ce{NH3}:
\begin{tightcenter}
    \ce{N2H4(aq) + H2O(l) <=> N2H5+(aq) + OH-(aq)}\ \ \ $K_{\text b1}=8.5\times10^{-7}$\\
    \ce{N2H5+(aq) + H2O(l) <=> N2H6^2+(aq) + OH-(aq)}\ \ \ $K_{\text b1}=8.9\times10^{-16}$
\end{tightcenter}
这种现象与~\ce{H2O}, ~\ce{H2O2}的酸性强弱关系类似.
\paragraph{联氨的制备与反应}
一般可以通过下面的方法制备~\ce{N2H4}:
\begin{tightcenter}
    \ce{2NH3 + NaOCl ->T[强碱水溶液] N2H4 + NaCl + H2O}
\end{tightcenter}
这一反应可能的机理如下:
\begin{tightcenter}
    \ce{NH3 + ClO- -> NH2Cl + OH-}\\
    \ce{NH3 + NH2Cl -> N2H5+ + Cl-}\\
    \ce{N2H5+ + OH- -> N2H4 + H2O}
\end{tightcenter}
联氨可以通过如下方式被滴定:
\begin{tightcenter}
    \ce{N2H4 + KIO3 + 2HCl ->T[\ce{H2O}/\ce{CCl4}] N2 + KCl + ICl + 3H2O}
\end{tightcenter}
以~\ce{KIO3}作为滴定剂,反应首先生成~\ce{I2}使得~\ce{CCl4}相变为紫色,然后继续加入~\ce{KIO3}又将其氧化为~\ce{ICl},利用~\ce{CCl4}相中~\ce{I2}颜色的完全消失判断终点.
\subsubsection{羟胺~\ce{NH2OH}}
由于羟胺的性质与联氨有一定相似之处,因此放在氢化物一节一同介绍.羟胺~\ce{NH2OH}是对热不稳定的白色固体,熔点为$\SI{33}{\celsius}33$,沸点$\SI{110}{\celsius}33$.
\paragraph{羟胺的结构}
\ce{NH2OH}的结构在在不同情况下差异较大.在晶体中,~\ce{NH2OH}倾向于以反式的构象.
\paragraph{羟胺的碱性}
羟胺的碱性比\ce{N2H4}或\ce{NH3}都要弱:
\begin{tightcenter}
    \ce{NH2OH(aq) + H2O(l) <=> NH3OH+(aq) + OH-(aq)}\ \ \ $K_\text b=6.6\times10^{-9}$
\end{tightcenter}
这是因为~\ce{O}有比~\ce{N}更强的吸电子效应.
\paragraph{羟胺的制备与反应}
\ce{NH2OH}可以通过许多反应制得.例如Raschig法合成~\ce{NH2OH}:
\begin{tightcenter}
    \ce{NH4NO2 + 2SO2 + NH3 + H2O -> (NH4)2[N(OH)(OSO2)2]}\\
    \ce{(NH4)2[N(OH)(OSO2)2] + 2H2O -> NH2OH + 2NH4HSO4}
\end{tightcenter}
或者利用下面几种方法制备羟胺盐:
\begin{tightcenter}
    \ce{2NO(g) + 3H2(g) + H2SO4(aq) ->T[\ce{Pt/C}] (NH3OH)2SO4} \\
    \ce{2RCH2NO2 + 2H2O + H2SO4 -> 2RCOOH + (NH3OH)2SO4}\\
    \ce{C2H4 + 2NO2 -> O2NCH2CH2NO2}\ \ \ \ \ \ce{O2NCH2CH2NO2 + H2SO4 -> CO2 + CO + (NH3OH)2SO4}
\end{tightcenter}
实验室中常用~\ce{HSO3-}还原~\ce{HNO2}制备,这与Raschig法类似.
\subsubsection{叠氮酸\ce{HN3}及其盐}
叠氮酸~\ce{HN3}是一种具有令人厌恶的刺激性气味的液体或气体,熔点$\SI{-80}{\celsius}$.~\ce{HN3}具有强烈的爆炸性,并且是一种剧毒的物质.\\
\indent 最常见的叠氮盐是叠氮化钠~\ce{NaN3},是一种剧毒且易爆的白色固体,加热至$\SI{275}{\celsius}$时分解.
\paragraph{\ce{HN3}与~\ce{N3-}的结构}
气相中的~\ce{HN3}为折线形分子,~\ce{N}原子在同一直线上,~\ce{HN-N2}键长$\SI{124}{\pico\meter}$,\ce{HN2-N}键长$\SI{113}{\pico\meter}$.键角为$112^\circ$.\\
\indent \ce{N3-}作为~\ce{CO2}的等电子体是直线型离子,其中~\ce{N-N}键键长均为$\SI{116}{\pico\meter}$.
\paragraph{\ce{HN3}的酸性}
~\ce{HN3}在水溶液中为弱酸,酸性与醋酸比较接近:
\begin{tightcenter}
    \ce{HN3(aq) + H2O(aq) -> H3O+(aq) + N3-(aq)}\quad$K_{a}=1.8\times10^5$,$\mathrm{p}K_{a}=4.77$
\end{tightcenter}
\paragraph{\ce{HN3}与~\ce{N3-}的制备与性质}
叠氮酸最早被通过以下方法制得:
\begin{tightcenter}
    \ce{N2H5+ + HNO2 -> HN3 + H+ +2H2O}
\end{tightcenter}
不过通常可以通过小心酸化~\ce{NaN3}的方式制备.大多数制取叠氮酸及其衍生物的方法都要用到~\ce{NaN3},它可以通过如下方式制备:
\begin{tightcenter}
    \ce{NaNO2 + 3NaNH2 ->T[熔融] NaN3 + 3NaOH + NH3}\\
    \ce{N2O + 2NaNH2 ->T[熔融] NaN3 + NaOH + NH3}(在\ce{NH3(l)}中进行反应亦可)\\
    \ce{3N2O + 4Na + NH3 -> NaN3 + 3NaOH + 2N2}
\end{tightcenter}
叠氮化物的主要用途取决于其爆炸性.例如,汽车安全气囊中曾经就置有~\ce{NaN3},分解时放出大量~\ce{N2}充满气囊,起到缓冲作用:
\begin{tightcenter}
    \ce{2NaN3 -> 2Na + 3N2}
\end{tightcenter}
而~\ce{Pb(NO3)2}与~\ce{NaN3}的复分解反应的产物~\ce{Pb(N3)2}则由于其可靠性(尤其是在潮湿环境中)而被广泛地用作引爆剂.
\subsubsection{二亚胺~\ce{N2H2}}
二亚胺~\ce{N2H2}是一种极不稳定的氮氢化物.目前,只有反式的~\ce{N2H2}被分离得到,它是一种无色的气体,在室温下即迅速分解,分解反应的半衰期只有数分钟.
\paragraph{二亚胺的研究历史}
1892年, Thiele在研究偶氮二甲酸的分解时发现最终产物只有~\ce{N2H4},~\ce{N2}和~\ce{CO2}\cite{Thiele1892}:
\begin{tightcenter}
    \ce{2HOOC-N=N-COOH -> N2H4 + N2 + 4CO2}
\end{tightcenter}
他因此认为反应经历了~\ce{N2H2}中间体,然后迅速歧化为~\ce{N2H4}和~\ce{N2}.此后,又有一些研究间接表明~\ce{N2H2}的存在,但直到1958年, Foner和Hudson才在~\ce{NH3}或~\ce{N2H4}的放电分解产物中检测到了~\ce{N2H2}\cite{Foner1958}.\\
\indent 1961年, Corey等人发现溶液中原位产生的~\ce{N2H2}可以高效地为烯烃顺式加氢\cite{Corey1961},立体选择性来源于反应协同进行.二亚胺还原自此成为有机化学中的一个独立的反应,并被广泛应用于烯烃的选择性还原中.
\paragraph{二亚胺的结构}
二亚胺有顺式和反式两种结构和一种可能的异构体:
\begin{figure}[H]\centering
    \begin{subfigure}[c]{.32\linewidth}
        \includegraphics{figure/N2H2-1.eps}\caption{反式~\ce{N2H2}}
    \end{subfigure}
    \begin{subfigure}[c]{.32\linewidth}
        \includegraphics{figure/N2H2-2.eps}\caption{顺式~\ce{N2H2}}
    \end{subfigure}
    \begin{subfigure}[c]{.32\linewidth}
        \includegraphics{figure/N2H2-3.eps}\caption{\ce{N2H2}的可能的异构体}
    \end{subfigure}
    \caption{~\ce{N2H2}的三种结构}
\end{figure}
在上述三种结构中最稳定的是反式~\ce{N2H2},因为其中存在两组~\ce{N}的孤对电子对$\sigma^\ast(\ce{N-H})$的超共轭作用.然而,在前面讲到的二亚胺还原中参与反应的却是顺式的~\ce{N2H2},因为顺式~\ce{N2H2}对烯烃的还原是协同进行的,具有更低的活化能.在质子溶剂中,二者通过可逆的酸碱反应发生转化,反应性更高的顺式~\ce{N2H2}不断被消耗,而反式~\ce{N2H2}不断转化为顺式~\ce{N2H2}.
\end{document}